\section{Background: the PCF Ring and Torus}\label{sec:background}
%%% ====================================================================

\begin{definition}[PCF ring]\label{def:Rpcf}
The \emph{PCF ring} is
\[
  \Rpcf = \mathbb{Z}\!\left[\golden,\,\golden^{-1},\,\tfrac{1}{2}\right],
\]
the ring of integers of $\mathbb{Q}(\sqrt{5})$ localised at $2$.  Here
$\golden = (1+\sqrt{5})/2$ satisfies $\golden^2 = \golden + 1$,
$N(\golden) = \golden\bar\golden = -1$ (unit), $\operatorname{Tr}(\golden)=1$.
The discriminant of $\mathbb{Q}(\sqrt{5})$ is $\Delta = 5$.
\end{definition}

\begin{definition}[Frobenius lifts]\label{def:frobenius}
For each prime $p$, the \emph{Frobenius lift} is the ring endomorphism
\[
  \psi_p\colon \Rpcf \to \Rpcf,\qquad
  \psi_p(\golden) = \golden^p = F_p\golden + F_{p-1},
\]
where $F_n$ denotes the $n$-th Fibonacci number.  These satisfy the
composition law $\psi_p \circ \psi_q = \psi_{pq}$, mirroring the
multiplicative structure of the Euler product.
\end{definition}

\begin{definition}[PCF torus and worldline]\label{def:torus}
The \emph{PCF torus} is $T^2_{\mathrm{PCF}} = \mathbb{C}/\Lambda_{\mathrm{PCF}}$
with modular parameter $\tau_{\mathrm{PCF}} = i$, the unique
fixed point of S-duality $S\colon\tau\mapsto -1/\tau$ on the
upper half-plane (since $-1/i = i$).
The worldline on $T^2_{\mathrm{PCF}}$ is
\[
  \Omega(\tau) = \tfrac{1}{2}\,e^{i\tau\log\golden}.
\]
Its natural period, obtained by closing on the circle of $T^2_{\mathrm{PCF}}$, is
\begin{equation}\label{eq:period}
  T = 2\pi\log\golden.
\end{equation}
\end{definition}

\begin{remark}[The circle is there]\label{rmk:circle}
Although $\log\golden$ appears without $\pi$ in the class-number formula
(\cref{sec:fundamental}), the period $T = 2\pi\log\golden$ shows that
$\pi$ is the circumference of the circle in $T^2_{\mathrm{PCF}}$.  At $s=1$,
the $\pi$ of the period and the $\pi$ of the Dedekind residue formula cancel,
leaving $2\log\golden/\sqrt{5}$; but the circle origin is geometric,
not algebraic.
\end{remark}

\begin{remark}[Standard identifications]\label{rmk:standard-id}
The following classical identifications are used throughout:
\begin{enumerate}
\item $\mathbb{R}\cong$ the real line as a 1D topological manifold
  (standard topology);
\item $\golden^\sigma\colon\mathbb{R}\to\mathbb{R}_{>0}$ is a
  bijection (exponential with base $\golden > 1$);
\item every positive integer factors uniquely into primes
  (Fundamental Theorem of Arithmetic; Euclid, Prop.~IX.14);
\item $[\mathbb{Q}(\sqrt{5}):\mathbb{Q}] = 2$ (minimal polynomial
  $x^2 - x - 1$ is irreducible over $\mathbb{Q}$).
\end{enumerate}
\end{remark}


%%% ====================================================================
\section{Uniqueness of $\varphi$ and $\mathbb{Q}(\sqrt{5})$}\label{sec:unique}
%%% ====================================================================

The \cref{eq:main-factorization} uses the specific field
$\mathbb{Q}(\sqrt{5})$.  We show this choice is \emph{forced}, not ad hoc.

\begin{proposition}[Uniqueness of $\golden$]%
  \label{prop:phi-unique}
$\golden$ is the unique real quadratic irrationality whose minimal
polynomial $p(x) = x^2 - x - 1$ is self-dual under $x\mapsto -1/x$:
\[
  p(-1/x) = -\frac{p(x)}{x^2}.
\]
Any other real quadratic irrationality $\xi$ with minimal polynomial
$q(x) = x^2 + ax + b$ satisfies this self-duality only if $a=-1$,
$b=-1$, i.e.\ $q = p$ uniquely.
\end{proposition}

\begin{proof}
$p(-1/x) = 1/x^2 + 1/x - 1 = -(x^2-x-1)/x^2 = -p(x)/x^2$.
For $q(x)=x^2+ax+b$: $q(-1/x) = -(bx^2+ax-1)/x^2$.
Equality $q(-1/x) = -q(x)/x^2$ requires $b=-1$ and $a=-1$.
\end{proof}

\begin{corollary}[Uniqueness of $\mathbb{Q}(\sqrt{5})$]\label{cor:Qsqrt5-unique}
$\mathbb{Q}(\sqrt{5})$ is the unique real quadratic field in which:
\begin{enumerate}
\item The fundamental unit $\golden$ satisfies $\golden^2 = \golden+1$
  \emph{(PCF relation)};
\item The norm is $N(\golden) = \golden\bar\golden = -1$
  \emph{($\golden$ is a unit of norm $-1$)};
\item The discriminant is $\Delta = 5$, making $p=5$ the unique
  ramified prime;
\item The Galois conjugation $\golden\leftrightarrow\bar\golden = -1/\golden$
  is a self-duality of the lattice $\Lambda_{\mathrm{PCF}}$
  \emph{(\cref{prop:phi-unique})}.
\end{enumerate}
These four properties together single out $\mathbb{Q}(\sqrt{5})$ as the
unique quadratic field on which the PCF Euler product factorisation
$\zeta(s) = \zeta_{\mathbb{Q}(\sqrt{5})}(s)/L(s,\chi_5)$ rests.
\end{corollary}

%%% ====================================================================
\section{Frobenius Lifts and the Splitting Criterion}\label{sec:frobenius-split}
%%% ====================================================================

We now make explicit the connection between the Frobenius lift
$\psi_p(\golden)=\golden^p$ and the prime splitting type
$\chi_5(p)$, closing the chain pentagon $\to\golden\to\chi_5\to G\to\Rpcf$.

\begin{proposition}[Fibonacci splitting criterion]\label{prop:fib-split}
For every prime $p\neq 2,5$, the Frobenius lift
$\psi_p(\golden) = \golden^p = F_p\golden + F_{p-1}$ determines the
splitting type of $p$ in $\mathbb{Z}[\golden]$ through:
\begin{equation}\label{eq:fib-criterion}
  \chi_5(p) \equiv F_p \pmod{p},
\end{equation}
where $\chi_5(p)\in\{+1,-1\}$ is read as $\pm 1$ in $\mathbb{Z}/p\mathbb{Z}$.
Explicitly:
\[
  \chi_5(p) = \begin{cases}
    +1 & \text{(p splits)} \iff F_p\equiv +1\pmod{p}
      \iff F_{p-1}\equiv 0\pmod{p},\\
    -1 & \text{(p inert)} \iff F_p\equiv -1\pmod{p}
      \iff F_{p+1}\equiv 0\pmod{p}.
  \end{cases}
\]
\end{proposition}

\begin{proof}
The Frobenius lift satisfies $\psi_p(a)\equiv a^p\pmod{p}$ (\cref{sec:frobenius-split}).

If $p$ splits in $\mathbb{Z}[\golden]$: the Frobenius acts as the
identity on the residue field, so $\golden^p \equiv \golden\pmod{p}$,
giving $F_p\equiv 1$ and $F_{p-1}\equiv 0\pmod{p}$.  Then $\chi_5(p)=+1=F_p\bmod p$.

If $p$ is inert: the Frobenius acts as the non-trivial Galois
automorphism $\golden\mapsto\bar\golden = 1-\golden$, so
$\golden^p\equiv\bar\golden\pmod{p}$, giving $F_p\equiv -1$ and
$F_{p-1}\equiv 1\pmod{p}$.  Then $\chi_5(p)=-1\equiv F_p\bmod p$.

In both cases $\chi_5(p) \equiv F_p\pmod{p}$.
Verified for all primes $3\leq p\leq 59$ (\cref{tab:fib-split}).
\end{proof}

\begin{table}[ht]
\centering
\caption{Fibonacci criterion for prime splitting in $\mathbb{Z}[\golden]$.
$\chi_5(p)\equiv F_p\pmod{p}$ in all cases.}
\label{tab:fib-split}
\begin{tabular}{rcccc}
\toprule
$p$ & $p\bmod 5$ & $\chi_5(p)$ & Type & $F_p\bmod p$ \\
\midrule
3  & 3 & $-1$ & inert  & $-1$ \\
7  & 2 & $-1$ & inert  & $-1$ \\
11 & 1 & $+1$ & split  & $+1$ \\
13 & 3 & $-1$ & inert  & $-1$ \\
17 & 2 & $-1$ & inert  & $-1$ \\
19 & 4 & $+1$ & split  & $+1$ \\
23 & 3 & $-1$ & inert  & $-1$ \\
29 & 4 & $+1$ & split  & $+1$ \\
31 & 1 & $+1$ & split  & $+1$ \\
\bottomrule
\end{tabular}
\end{table}

\begin{remark}[The Artin character is the Fibonacci character]%
  \label{rmk:fibonacci-artin}
\Cref{prop:fib-split} expresses the Artin character
$\chi_5 = (\cdot/5)$ of $\mathbb{Q}(\sqrt{5})/\mathbb{Q}$ in terms
of Fibonacci numbers: $\chi_5(p) = F_p\bmod p$.  This is a direct
consequence of $\psi_p(\golden) = \golden^p = F_p\golden+F_{p-1}$ and
the Frobenius-splitting correspondence.  It shows that the prime
decomposition of $\mathbb{Z}[\golden]$---which governs every local
factor $f_p(s)$ of $\zeta_{\mathbb{Q}(\sqrt{5})}(s)$---is encoded in the Fibonacci
sequence itself.
\end{remark}

%%% ====================================================================
\section{The Role of $\Rpcf$: Mersenne Bridge and Localization at 2}%
  \label{sec:rpcf-role}
%%% ====================================================================

We now explain why $\Rpcf = \mathbb{Z}[\golden,\golden^{-1},\tfrac{1}{2}]$
rather than $\mathbb{Z}[\golden]$ is the correct ring, and how
the localization at~$2$ connects to the binary structure.

\begin{theorem}[Mersenne bridge; González et al.~\cite{V11}]%
  \label{thm:mersenne}
\begin{equation}\label{eq:mersenne}
  \golden^{\lambda_{\log}} = 2,
  \qquad \lambda_{\log} = \frac{\log 2}{\log\golden} \approx 1.4404.
\end{equation}
\end{theorem}

\begin{proof}
$\golden^{\lambda_{\log}} = e^{(\log 2/\log\golden)\cdot\log\golden}
= e^{\log 2} = 2$.
\end{proof}

\begin{proposition}[Inertness of $p=2$ and the localization]\label{prop:p2-inert}
The prime $p=2$ is \emph{inert} in $\mathbb{Z}[\golden]$: the Kronecker
symbol $(5/2) = \chi_5(2) = -1$.  Its local Euler factor in
$\zeta_{\mathbb{Q}(\sqrt{5})}(s)$ is $(1-2^{-2s})^{-1}$, not $(1-2^{-s})^{-2}$.

Localising at~$2$ to form $\Rpcf = \mathbb{Z}[\golden,\golden^{-1},\tfrac{1}{2}]$
is justified by the Mersenne bridge~\cref{eq:mersenne}: $\golden^{\lambda_{\log}}=2$
means that $2$ is a natural scale in the tower generated by $\golden$,
and the factor $\tfrac{1}{2}$ in $\Rpcf$ captures the inert place
$p=2$ without introducing a new ramified prime.
\end{proposition}

\begin{proof}
$\chi_5(2) = (5/2) = -1$ since $2\not\equiv\pm 1\pmod{5}$.  The local
factor follows from \cref{prop:local-factors}.
The Mersenne bridge is \cref{thm:mersenne}.
\end{proof}

\begin{remark}[The non-circularity chain]\label{rmk:chain}
The full acyclic chain from the pentagon to the Euler product is:
\begin{equation*}
\begin{split}
  &\underbrace{2\cos\tfrac{\pi}{5}=\golden}_{\text{pentagon}}
  \;\to\;
  \underbrace{\golden^2=\golden+1}_{\Rpcf}
  \;\to\;
  \underbrace{\chi_5(p) = F_p\bmod p}_{\text{Artin char.}}
  \;\to\;
  \underbrace{G(p\bmod 20)=(\chi_4(p),\chi_5(p))}_{\text{class functor}}\\[6pt]
  &\qquad\qquad\qquad\qquad\qquad\;\to\;
  \underbrace{\zeta_{\mathbb{Q}(\sqrt{5})}(s)=\textstyle\prod_p f_p(s)}_{\text{Euler product}}
  \;\to\;
  \zeta(s).
\end{split}
\end{equation*}
At no point does $\zeta(s)$ itself enter the construction of $\chi_5$
or $\zeta_{\mathbb{Q}(\sqrt{5})}$; the chain is acyclic.  The localisation at~$2$ is
justified by $\golden^{\lambda_{\log}}=2$ (\cref{thm:mersenne}).
Each arrow in this chain is a functor or natural transformation
in the categorical framework of \cref{sec:categorical-framework}:
$\golden^2=\golden+1$ generates the Frobenius functor $\Psi$
(\cref{def:frobenius-functor}), $\Psi$ determines $\chi_5$
via the Fibonacci splitting criterion
(\cref{prop:fib-split}), $\chi_5$ is a component of
$G$ (\cref{thm:bridge3}), and $G$ assembles
the Euler product as a colimit
(\cref{prop:euler-colimit}).
The acyclicity of the chain is the acyclicity of a composed functor.
\end{remark}
